\uline{is an isomorphism}.

To see that it is injective, we use the sequence of canonical homomorphisms:
@@M0@@
the last arrow coming from the fact that @@M1@@ is a subfunctor of @@M2@@. To verify that the composite homomorphism is the identity, we observe that a morphism of functors: @@M3@@ is determined by its value on the objects of @@M4@@ whose filtration is coarse; indeed, let @@M5@@, and let @@M6@@ be the algebra @@M7@@ endowed with the coarse filtration; the canonical morphism @@M8@@ defines a commutative diagram:
@@M9@@
showing that @@M10@@ is the restriction of @@M11@@ to @@M12@@.

We next consider the composite homomorphism:
@@M13@@
and it is clear that this is the identity, giving the asserted result.

\sgaStableTarget{sga6.e05.thm.1-10}\textbf{Theorem 1.10.} \uline{Let} @@M14@@ \uline{be a functor on the category} @@M15@@ \uline{of} @@M16@@-\uline{algebras}, \uline{with values in the category of separated complete filtered abelian groups}, \uline{and commuting with invariants under a finite group}. \uline{Then the map from} @@M17@@ \uline{to} @@M18@@ \uline{which}, \uline{to a morphism of functors}, \uline{assigns}
